\section{AIME II 2013}
        \begin{enumerate}
        	\item (\textit{AIME II 2013}) Suppose that the measurement of time during the day is converted to the metric system so that each day has $10$ metric hours, and each metric hour has $100$ metric minutes. Digital clocks would then be produced that would read $\text{9:99}$ just before midnight, $\text{0:00}$ at midnight, $\text{1:25}$ at the former $\text{3:00}$ AM, and $\text{7:50}$ at the former $\text{6:00}$ PM. After the conversion, a person who wanted to wake up at the equivalent of the former $\text{6:36}$ AM would set his new digital alarm clock for $\text{A:BC}$, where $\text{A}$, $\text{B}$, and $\text{C}$ are digits. Find $100\text{A}+10\text{B}+\text{C}$.


 
 \textbf{Answer: }275



\textbf{Solution}

There are $24 \cdot 60=1440$ normal minutes in a day , and $10 \cdot 100=1000$ metric minutes in a day. The ratio of normal to metric minutes in a day is $\frac{1440}{1000}$, which simplifies to $\frac{36}{25}$. This means that every time 36 normal minutes pass, 25 metric minutes pass. From midnight to $\text{6:36}$ AM, $6 \cdot 60+36=396$ normal minutes pass. This can be viewed as $\frac{396}{36}=11$ cycles of 36 normal minutes, so 11 cycles of 25 metric minutes pass. Adding $25 \cdot 11=275$ to $\text{0:00}$ gives $\text{2:75}$, so the answer is $\boxed{275}$.

 





\textbf{Solution 2}

First we want to find out what fraction of a day has passed at 6:36 AM. One hour is $\frac{1}{24}$ of a day, and 36 minutes is $\frac{36}{60}=\frac{3}{5}$ of an hour, so at 6:36 AM, $6 \cdot \frac{1}{24} + \frac{1}{24} \cdot \frac{3}{5}=\frac{1}{4}+\frac{1}{40}=\frac{11}{40}$ of a day has passed. Now the metric timing equivalent of $\frac{11}{40}$ of a day is $\frac{11}{40}\cdot 1000=275$ metric minutes, which is equivalent to 2 metric hours and 75 metric minutes, so our answer is $\boxed{275}$ - mathleticguyyy

 




	\item (\textit{AIME II 2013}) Positive integers $a$ and $b$ satisfy the condition

 \[\log_2(\log_{2^a}(\log_{2^b}(2^{1000}))) = 0.\]
Find the sum of all possible values of $a+b$.


 
 \textbf{Answer: }881



\textbf{Solution 1}

To simplify, we write this logarithmic expression as an exponential one. Just looking at  the first log, it has a base of 2 and an argument of the expression in parenthesis. Therefore, we can make 2 the base, 0 the exponent, and the argument the result. That means $\log_{2^a}(\log_{2^b}(2^{1000}))=1$ (because $2^0=1$). Doing this again, we get $\log_{2^b}(2^{1000})=2^a$. Doing the process one more time, we finally eliminate all of the logs, getting ${(2^{b})}^{(2^a)}=2^{1000}$. Using the property that ${(a^x)^{y}}=a^{xy}$, we simplify to $2^{b\cdot2^{a}}=2^{1000}$. Eliminating equal bases leaves $b\cdot2^a=1000$. The largest $a$ such that $2^a$ divides $1000$ is $3$, so we only need to check $1$,$2$, and $3$. When $a=1$, $b=500$; when $a=2$, $b=250$; when $a=3$, $b=125$. Summing all the $a$'s and $b$'s gives the answer of $\boxed{881}$.

 Note that $a$ cannot be $0,$ since that would cause the $\log_{2^a}$ to have a $1$ in the base, which is not possible (also the problem specifies that $a$ and $b$ are positive).

 





\textbf{Solution 2}

We proceed as in Solution 1, raising $2$ to both sides to achieve $\log_{2^a}(\log_{2^b}(2^{1000})) = 1.$ We raise $2^a$ to both sides to get $\log_{2^b}(2^{1000})=2^a$, then simplify to get $\dfrac{1000}b=2^a$. 

 At this point, we want both $a$ and $b$ to be integers. Thus, $2^a$ can only be a power of $2$. To help us see the next step, we factorize $1000$: $\dfrac{2^35^3}b=2^a.$ It should be clear that $a$ must be from $1$ to $3$; when $a=1$, $b=500$; when $a=2$, $b=250$; and finally, when $a=3$, $b=125.$ We sum all the pairs to get $\boxed{881}.$ 

 




	\item (\textit{AIME II 2013}) A large candle is $119$ centimeters tall.  It is designed to burn down more quickly when it is first lit and more slowly as it approaches its bottom.  Specifically, the candle takes $10$ seconds to burn down the first centimeter from the top, $20$ seconds to burn down the second centimeter, and $10k$ seconds to burn down the $k$-th centimeter.  Suppose it takes $T$ seconds for the candle to burn down completely.  Then $\frac{T}{2}$ seconds after it is lit, the candle's height in centimeters will be $h$.  Find $10h$.


 
 \textbf{Answer: }350



\textbf{Solution}

We find that $T=10(1+2+\cdots +119)$. From Gauss's formula, we find that the value of $T$ is $10(7140)=71400$. The value of $\frac{T}{2}$ is therefore $35700$. We find that $35700$ is $10(3570)=10\cdot \frac{k(k+1)}{2}$, so $3570=\frac{k(k+1)}{2}$. As a result, $7140=k(k+1)$, which leads to $0=k^2+k-7140$. We notice that $k=84$, so the answer is $10(119-84)=\boxed{350}$.

 




	\item (\textit{AIME II 2013}) In the Cartesian plane let $A = (1,0)$ and $B = \left( 2, 2\sqrt{3} \right)$.  Equilateral triangle $ABC$ is constructed so that $C$ lies in the first quadrant.  Let $P=(x,y)$ be the center of $\triangle ABC$.  Then $x \cdot y$ can be written as $\frac{p\sqrt{q}}{r}$, where $p$ and $r$ are relatively prime positive integers and $q$ is an integer that is not divisible by the square of any prime.  Find $p+q+r$.


 
 \textbf{Answer: }040



\textbf{Solution 1}

The distance from point $A$ to point $B$ is $\sqrt{13}$. The vector that starts at point A and ends at point B is given by $B - A = (1, 2\sqrt{3})$. Since the center of an equilateral triangle, $P$, is also the intersection of the perpendicular bisectors of the sides of the triangle, we need first find the equation for the perpendicular bisector to $\overline{AB}$. The line perpendicular to $\overline{AB}$ through the midpoint, $M =  \left(\dfrac{3}{2},\sqrt{3}\right)$, $\overline{AB}$ can be parameterized by $\left(\dfrac{2\sqrt{3}}{\sqrt{13}}, \dfrac{-1}{\sqrt{13}}\right) t + \left(\dfrac{3}{2},\sqrt{3}\right)$. At this point, it is useful to note that $\Delta BMP$ is a 30-60-90 triangle with $\overline{MB}$ measuring $\dfrac{\sqrt{13}}{2}$. This yields the length of $\overline{MP}$ to be $\dfrac{\sqrt{13}}{2\sqrt{3}}$. Therefore, $P =\left(\dfrac{2\sqrt{3}}{\sqrt{13}},\dfrac{-1}{\sqrt{13}}\right)\left(\dfrac{\sqrt{13}}{2\sqrt{3}}\right) + \left(\dfrac{3}{2},\sqrt{3}\right) = \left(\dfrac{5}{2}, \dfrac{5}{2\sqrt{3}}\right)$. Therefore $xy = \dfrac{25\sqrt{3}}{12}$ yielding an answer of $p + q + r  = 25 + 3 + 12 = \boxed{040}$.

 





\textbf{Solution 2}

Rather than considering the Cartesian plane, we use complex numbers. Thus A is 1 and B is $2 + 2\sqrt{3}i$. 

 Recall that a rotation of $\theta$ radians counterclockwise is equivalent to multiplying a complex number by $e^{i\theta}$, but here we require a clockwise rotation, so we multiply by $e^{-\frac{i\pi}{3}}$ to obtain C. Upon averaging the coordinates of A, B, and C, we obtain the coordinates of P, viz. $\left(\frac{5}{2}, \frac{5\sqrt{3}}{6}\right)$.

 Therefore $xy$ is $\frac{25\sqrt{3}}{12}$ and the answer is $25 + 12 + 3 = \boxed{040}$.

 





\textbf{Solution 3}

We can also consider the slopes of the lines. Midpoint $M$ of $AB$ has coordinates $\left(\frac{3}{2},\ \sqrt{3}\right)$. Because line $AB$ has slope $2\sqrt{3}$, the slope of line $MP$ is  $-\frac{1}{2\sqrt{3}}$ (Because of perpendicular slopes). 

 Since $\Delta ABC$ is equilateral, and since point $P$ is the centroid, we can quickly calculate that $MP = \frac{\sqrt{39}}{6}$. Then, define $\Delta x$ and $\Delta y$ to be the differences between points $M$ and $P$. Because of the slope, it is clear that $\Delta x = 2\sqrt{3} \Delta y$.

 We can then use the Pythagorean Theorem on line segment $MP$: $MP^2 = \Delta x^2 + \Delta y^2$ yields $\Delta y = -\frac{1}{2\sqrt{3}}$ and $\Delta x = 1$, after substituting $\Delta x$. The coordinates of P are thus $\left(\frac{5}{2},\ \frac{5\sqrt{3}}{6}\right)$. Multiplying these together gives us $\frac{25\sqrt{3}}{12}$, giving us $\boxed{040}$ as our answer.

 





\textbf{Solution 4}

Since $AC$ will be segment $AB$ rotated clockwise $60^{\circ}$, we can use a rotation matrix to find $C$. We first translate the triangle $1$ unit to the left, so $A'$ lies on the origin, and $B' = (1, 2\sqrt{3})$. Rotating clockwise $60^{\circ}$ is the same as rotating $300^{\circ}$ counter-clockwise, so our rotation matrix is $\begin{bmatrix} \cos{300^{\circ}} & -\sin{300^{\circ}}\\ \sin{300^{\circ}} & \cos{300^{\circ}}\\ \end{bmatrix} = \begin{bmatrix} \frac{1}{2} & \frac{\sqrt{3}}{2}\\ -\frac{\sqrt{3}}{2} & \frac{1}{2}\\ \end{bmatrix}$. Then $C' = \begin{bmatrix} \frac{1}{2} & \frac{\sqrt{3}}{2}\\ -\frac{\sqrt{3}}{2} & \frac{1}{2}\\ \end{bmatrix} \cdot \begin{bmatrix} 1\\ 2\sqrt{3}\\ \end{bmatrix} = \begin{bmatrix} \frac{7}{2}\\ \frac{\sqrt{3}}{2}\\ \end{bmatrix}$. Thus, $C = (\frac{9}{2}, \frac{\sqrt{3}}{2})$. Since the triangle is equilateral, the center of the triangle is the average of the coordinates of the vertices. Then $P = (\frac{1 + 2 + \frac{9}{2}}{3}, \frac{2\sqrt{3} + \frac{\sqrt{3}}{2}}{3}) = (\frac{5}{2}, \frac{5\sqrt{3}}{6})$. Our answer is $\frac{5}{2} \cdot \frac{5\sqrt{3}}{6} = \frac{25\sqrt{3}}{12}$. $25 + 3 + 12 = \boxed{40}$

 





\textbf{Solution 5}

We construct point $C$ by drawing two circles with radius $r = AB = \sqrt{13}$.  One circle will be centered at $A$, while the other is centered at $B$.  The equations of the circles are:

 $(x - 1)^2 + y^2 = 13$

 $(x - 2)^2 + (y - 2\sqrt{3})^2 = 13$

 Setting the LHS of each of these equations equal to each other and solving for $x$ yields after simplification:

 $x = \frac{15}{2} - 2\sqrt{3}y$

 Plugging that into the first equation gives the following quadratic in $y$ after simplification:

 $y^2 - 2\sqrt{3}y + \frac{9}{4} = 0$

 The quadratic formula gives $y = \frac{\sqrt{3}}{2}, \frac{3\sqrt{3}}{2}$.  

 Since $x > 0$ and $x = \frac{15}{2} - 2\sqrt{3}y$, we pick $y = \frac{\sqrt{3}}{2}$ in the hopes that it will give $x > 0$.  Plugging $y$ into the equation for $x$ yields $x = \frac{9}{2}$.

 Thus, $C(\frac{9}{2}, \frac{\sqrt{3}}{2})$.  Averaging the coordinates of the vertices of equilateral triangle $ABC$ will give the center of mass of the triangle.  

 Thus, $P(\frac{5}{2}, \frac{5\sqrt{3}}{6})$, and the product of the coordinates is $\frac{25\sqrt{3}}{12}$, so the desired quantity is $\boxed{040}$. 

 





\textbf{Solution 6}

Labeling our points and sketching a graph we get that $C$ is to the right of $AB$. Of course, we need to find $C$. Note that the transformation from $A$ to $B$ is $[1,2\sqrt{3}]$, and if we imagine a height dropped to $AB$ we see that a transformation from the midpoint $(\frac{3}{2},\sqrt {3})$ to $C$ is basically the first transformation, with $\frac{\sqrt{3}}{2}$ the magnitude and the x and y switched– then multiply the new y by -1. Then, applying this transformation of $[3,\frac{-\sqrt{3}}{2}]$ we get that $C=(\frac{9}{2},\frac{\sqrt{3}}{2})$ which means that $P=(\frac{5}{2},\frac{5\sqrt{3}}{6})$. Then our answer is $\boxed{40}$.

 





\textbf{Solution 7}

Transform this into the complex plane and let $a=1, b=2+2\sqrt3 i$. We know that 3 complex numbers $a,b,c$ form an equilateral triangle if $a^2+b^2+c^2=ab+bc+ac$, so plugging in our values of $a,b$, we get $8\sqrt3 i - 7 +c^2 = 2+2\sqrt3 i + (3+2\sqrt 3i)c.$ Solving for $c$ using Wolfram Alpha, we find that the solutions are $c=\frac 92 + \frac{i\sqrt3}{2}, -\frac 32 + \frac{3i\sqrt3}{2}$. The first one is in the first quadrant, so $C\left( \frac 92, \frac{\sqrt3}{2} \right)$. The center is the average of the coordinates and we find that it is $\left(\frac{5}2, \frac{5\sqrt3}{6} \right)$. Then $xy = \frac{25\sqrt3}{12} \implies 40$.

 -bobthegod78, krwang, and Simplest14

 




	\item (\textit{AIME II 2013}) In equilateral $\triangle ABC$ let points $D$ and $E$ trisect $\overline{BC}$. Then $\sin(\angle DAE)$ can be expressed in the form $\frac{a\sqrt{b}}{c}$, where $a$ and $c$ are relatively prime positive integers, and $b$ is an integer that is not divisible by the square of any prime. Find $a+b+c$.


 
 \textbf{Answer: }020



\textbf{Solution 1}


                    \begin{figure}[H]
                        \centering
                        \includegraphics[width=0.5\textwidth]{Solution/AIME_2013_II_5_Solution_Figure_1.png}
                    \end{figure}
                    

 Without loss of generality, assume the triangle sides have length 3. Then the trisected side is partitioned into segments of length 1, making your computation easier. 

 Let $M$ be the midpoint of $\overline{DE}$. Then $\Delta MCA$ is a 30-60-90 triangle with $MC = \dfrac{3}{2}$, $AC = 3$ and $AM = \dfrac{3\sqrt{3}}{2}$. Since the triangle $\Delta AME$ is right, then we can find the length of $\overline{AE}$ by pythagorean theorem, $AE = \sqrt{7}$. Therefore, since $\Delta AME$ is a right triangle, we can easily find $\sin(\angle EAM) = \dfrac{1}{2\sqrt{7}}$ and $\cos(\angle EAM) = \sqrt{1-\sin(\angle EAM)^2}=\dfrac{3\sqrt{3}}{2\sqrt{7}}$. So we can use the double angle formula for sine, $\sin(\angle EAD) = 2\sin(\angle EAM)\cos(\angle EAM) = \dfrac{3\sqrt{3}}{14}$. Therefore, $a + b + c = \boxed{020}$.

 





\textbf{Solution 2}

We find that, as before, $AE = \sqrt{7}$, and also the area of $\Delta DAE$ is $\dfrac{1}{3}$ the area of $\Delta ABC$. Thus, using the area formula, $\dfrac{1}{2} \cdot 7 \cdot \sin(\angle EAD) = \dfrac{3\sqrt{3}}{4}$, and  $\sin(\angle EAD) = \dfrac{3\sqrt{3}}{14}$. Therefore, $a + b + c = \boxed{020}.$

 





\textbf{Solution 3}

Let A be the origin of the complex plane, B be $1+i\sqrt{3}$, and C be $2$. Also, WLOG, let D have a greater imaginary part than E. Then, D is $\frac{4}{3}+\frac{2i\sqrt{3}}{3}$ and E is $\frac{5}{3}+\frac{i\sqrt{3}}{3}$. Then, $\sin(\angle DAE) = Im\left(\dfrac{\frac{4}{3}+\frac{2i\sqrt{3}}{3}}{ \frac{5}{3}+\frac{i\sqrt{3}}{3}}\right) = Im\left(\frac{26+6i\sqrt{3}}{28}\right) =  \frac{3\sqrt{3}}{14}$. Therefore, $a + b + c = \boxed{020}$

 





\textbf{Solution 4}

Without loss of generality, say that the side length of triangle ABC is 3. EC is 1 and by the law of cosines, $AE^2=1+3^2-2(1)(3)\cos(\angle DAE)$ or $AE=\sqrt7$ The same goes for AD. DE equals 1 because AD and AE trisect BC. By the law of cosines, $\cos(\angle DAE)=(1-7-7)/-2(\sqrt7)(\sqrt7)=13/14$ Since $\sin^2(\angle DAE)=1-cos^2(\angle DAE)$ Then $\sin^2(\angle DAE)= 1-\frac{169}{196}=\frac{27}{196}$ So $\sin(\angle DAE)=\frac{3\sqrt{3}}{14}$. Therefore, $a + b + c = \boxed{020}$

 





\textbf{Solution 5 (Vectors)}

Setting up a convinient coordinate system, we let $A$ be at point $(0, 0)$, $B$ be at point $(3, 3\sqrt3)$, and $C$ be at point $(6, 0)$. Then $D$ and $E$ will be at points $(4, 2\sqrt3)$ and $(5, \sqrt3)$. Then $\cos(\angle DAE) = \frac{\vec{AD}\cdot\vec{AE}}{\|\vec{AD}\| \|\vec{AE}\|} = \frac{4\cdot5 + 2\sqrt{3}\cdot\sqrt{3}}{28}=\frac{13}{14}$. From here, we see that $\sin(\angle DAE) = \sqrt{1-\cos^2(\angle DAE)} = \frac{3\sqrt3}{14}\Longrightarrow\boxed{020}$

 





\textbf{Solution 6}

We first drop the altitude from $A$ to $BC$, and that evaluates to a length of $5$. Now, we can find $AD$ by Pythag. We have $AD = \sqrt{7}$. Now, we can set $\angle{DAE}$ to be $\theta$, and now we we can solve for $\sin{\theta}$. We will take advantage of the different ways to find the area of a triangle. We have the area of the triangle is $\dfrac{\dfrac{1}{2} \cdot \dfrac{3\sqrt3}{2}}{2} = \dfrac{3\sqrt3}{4}$. We can also express it as $\dfrac{1}{2} \cdot 7 \cdot \sin{\theta}$, and solving for $\sin{\theta}$, we get 
 \[\dfrac{3\sqrt3}{14} \Longrightarrow \boxed{020}\]

 




	\item (\textit{AIME II 2013}) Find the least positive integer $N$ such that the set of $1000$ consecutive integers beginning with $1000\cdot N$ contains no square of an integer.


 
 \textbf{Answer: }282



\textbf{Solution 1}

The difference between consecutive integral squares must be greater than 1000. $(x+1)^2-x^2\geq1000$, so $x\geq\frac{999}{2}\implies x\geq500$. $x=500$ does not work, so $x>500$. Let $n=x-500$. By inspection,  $n^2$ should end in a number close to but less than 1000 such that there exists $1000N$ within the difference of the two squares. Examine when $n^2=1000$. Then, $n=10\sqrt{10}$. One example way to estimate $\sqrt{10}$ follows. 

 \begin{verbatim}
$3^2=9$, so $10=(x+3)^2=x^2+6x+9$. $x^2$ is small, so $10=6x+9$. $x=1/6\implies \sqrt{10}\approx 19/6$. This is 3.16.
\end{verbatim}Then, $n\approx 31.6$. $n^2<1000$, so $n$ could be $31$. Add 500 to get the first square and 501 to get the second. Then, the two integral squares are $531^2$ and $532^2$. Checking, $531^2=281961$ and $532^2=283024$. $282,000$ straddles the two squares, which have a difference of 1063. The difference has been minimized, so $N$ is minimized $N=282000\implies\boxed{282}$

 





\textbf{Solution 2}

Let us first observe the difference between $x^2$ and $(x+1)^2$, for any arbitrary $x\ge 0$. $(x+1)^2-x^2=2x+1$. So that means for every $x\ge 0$, the difference between that square and the next square have a difference of $2x+1$. Now, we need to find an $x$ such that $2x+1\ge 1000$. Solving gives $x\ge \frac{999}{2}$, so $x\ge 500$. Now we need to find what range of numbers has to be square-free: $\overline{N000}\rightarrow \overline{N999}$ have to all be square-free.Let us first plug in a few values of $x$ to see if we can figure anything out. $x=500$, $x^2=250000$, and $(x+1)^2=251001$. Notice that this does not fit the criteria, because $250000$ is a square, whereas $\overline{N000}$ cannot be a square. This means, we must find a square, such that the last $3$ digits are close to $1000$, but not there, such as $961$ or $974$. Now, the best we can do is to keep on listing squares until we hit one that fits. We do not need to solve for each square: remember that the difference between consecutive squares are $2x+1$, so all we need to do is addition. After making a list, we find that $531^2=281961$, while $532^2=283024$. It skipped $282000$, so our answer is $\boxed{282}$.

 





\textbf{Solution 3}

Let $x$ be the number being squared. Based on the reasoning above, we know that $N$ must be at least $250$, so $x$ has to be at least $500$. Let $k$ be $x-500$. We can write $x^2$ as $(500+k)^2$, or $250000+1000k+k^2$. We can disregard $250000$ and $1000k$, since they won't affect the last three digits, which determines if there are any squares between $\overline{N000}\rightarrow \overline{N999}$. So we must find a square, $k^2$, such that it is under $1000$, but the next square is over $1000$. We find that $k=31$ gives $k^2=961$, and so $(k+1)^2=32^2=1024$. We can be sure that this skips a thousand because the $1000k$ increments it up $1000$ each time. Now we can solve for $x$: $(500+31)^2=281961$, while $(500+32)^2=283024$. We skipped $282000$, so the answer is $\boxed{282}$.

 





\textbf{Solution 4}

The goal is to find the least $N \in \mathbb{N}$ such that $\exists m \in \mathbb{N}$ where $m^2 + 1 \leq 1000N, 1000N + 1000 \leq (m+1)^2$.

 
Combining the two inequalities leads to $(m+1)^2 \geq m^2 + 1001, m \geq 500$.

 
Let $m = k + 500$, where $k \in \mathbb{W}$, then the inequalities become,

 $N \geq \frac{(k+500)^2 + 1}{1000} = \frac{k^2 + 1}{1000} + k + 250$, and

 $N \leq \frac{(k+501)^2}{1000} - 1 = \frac{(k+1)^2}{1000} + k + 250.$

 
For $k=31$, one can verify that $N = 282$ is the unique integer satisfying the inequalities.

 For $k \leq 30$, $k + 250 < \frac{k^2 + 1}{1000} + k + 250 \leq N$ $\leq \frac{(k+1)^2}{1000} + k + 250 \leq \frac{(30+1)^2}{1000} + k + 250 < k + 251$, 

 i.e., $k + 250 < N < k + 251$, a contradiction. 

 Note $k \geq 32$ leads to larger $N$(s).

 Hence, the answer is $\boxed{282}$.

 


 




	\item (\textit{AIME II 2013}) A group of clerks is assigned the task of sorting $1775$ files. Each clerk sorts at a constant rate of $30$ files per hour. At the end of the first hour, some of the clerks are reassigned to another task; at the end of the second hour, the same number of the remaining clerks are also reassigned to another task, and a similar assignment occurs at the end of the third hour. The group finishes the sorting in $3$ hours and $10$ minutes. Find the number of files sorted during the first one and a half hours of sorting.


 
 \textbf{Answer: }945



\textbf{Solution}

There are $x$ clerks at the beginning, and $t$ clerks are reassigned to another task at the end of each hour. So, $30x+30(x-t)+30(x-2t)+30\cdot\frac{10}{60} \cdot (x-3t)=1775$, and simplify that we get $19x-21t=355$.Now the problem is to find a reasonable integer solution. Now we know $x= \frac{355+21t}{19}$, so $19$ divides $355+21t$, AND as long as $t$ is a integer, $19$ must divide $2t+355$. Now, we suppose that $19m=2t+355$, similarly we get $t=\frac{19m-355}{2}$, and so in order to get a minimum integer solution for $t$, it is obvious that $m=19$ works. So we get $t=3$ and $x=22$.  One and a half hour's work should be $30x+15(x-t)$, so the answer is $\boxed{945}$.

 





\textbf{Solution 2}

We start with the same approach as solution 1 to get $19x-21t=355$. Then notice that $21t + 355 \equiv 0 \pmod{19}$, or $2t-6 \equiv 0 \pmod{19}$, giving the smallest solution at $t=3$. We find that $x=22$. Then the number of files they sorted will be $30x+15(x-t)=660+285=\boxed{945}.$

 





\textbf{Solution 3 (More rigorous explanation)}

From solution one, we can obtain $19x-21t=355$ where $x$ is the number of clerks at the beginning and $t$ is the amount of clerks removed at the end of each hour. Since the number of clerks removed must be less than the amount of clerks at the beginning, we can form the inequality $x > 3t$. From here, we can replace $21t$ with $7x$ and form the inequality $19x-7x<355$. This means that $x < \frac{355}{12}$ or $x \le 29$. We can rewrite our equation at the beginning to become $19x = 355+21t$. Because $355 \equiv 13 \pmod {19}$ and $21 \equiv 2 \pmod {19}$, we can figure out that the minimum value of $t$ will be 3 as $13 + 2*3 = 19$. From here, we can find $x = 22$ and solve that the number of files the clerks sorted in the first 1 hour and 30 minutes is $30(22) + 15(22 - 3) = 660 + 285 = \boxed{945}$. If we increase $t$ by $19$, we will have to increase $x$ by $21$ which results in $x$ being $43$, thus breaking our inequality of $x \le 29$. Therefore, the only solutions for x and t where both are integers are 22 and 3, respectively.

 




	\item (\textit{AIME II 2013}) A hexagon that is inscribed in a circle has side lengths $22$, $22$, $20$, $22$, $22$, and $20$ in that order. The radius of the circle can be written as $p+\sqrt{q}$, where $p$ and $q$ are positive integers. Find $p+q$.


 
 \textbf{Answer: }272



\textbf{Solution}


                    \begin{figure}[H]
                        \centering
                        \includegraphics[width=0.5\textwidth]{Solution/AIME_2013_II_8_Solution_Figure_1.png}
                    \end{figure}
                    

 


 





\textbf{Solution 1}

Let us call the hexagon $ABCDEF$, where $AB=CD=DE=AF=22$, and $BC=EF=20$. We can just consider one half of the hexagon, $ABCD$, to make matters simpler. Draw a line from the center of the circle, $O$, to the midpoint of $BC$, $X$. Now, draw a line from $O$ to the midpoint of $AB$, $Y$. Clearly, $\angle BXO=90^{\circ}$, because $BO=CO$, and $\angle BYO=90^{\circ}$, for similar reasons. Also notice that $\angle AOX=90^{\circ}$. Let us call $\angle BOY=\theta$. Therefore, $\angle AOB=2\theta$, and so $\angle BOX=90-2\theta$. Let us label the radius of the circle $r$. This means 
 \[\sin{\theta}=\frac{BY}{r}=\frac{11}{r}\] 
 \[\sin{(90-2\theta)}=\frac{BX}{r}=\frac{10}{r}\]Now we can use simple trigonometry to solve for $r$.Recall that  $\sin{(90-\alpha)}=\cos(\alpha)$: That means $\sin{(90-2\theta)}=\cos{2\theta}=\frac{10}{r}$.Recall that $\cos{2\alpha}=1-2\sin^2{\alpha}$: That means $\cos{2\theta}=1-2\sin^2{\theta}=\frac{10}{r}$.Let $\sin{\theta}=x$.Substitute to get $x=\frac{11}{r}$ and $1-2x^2=\frac{10}{r}$Now substitute the first equation into the second equation: $1-2\left(\frac{11}{r}\right)^2=\frac{10}{r}$Multiplying both sides by $r^2$ and reordering gives us the quadratic 
 \[r^2-10r-242=0\]Using the quadratic equation to solve, we get that $r=5+\sqrt{267}$ (because $5-\sqrt{267}$ gives a negative value), so the answer is $5+267=\boxed{272}$.

 





\textbf{Solution 2 (Extremely Easy Geo)}

Using the trapezoid $ABCD$ mentioned above, draw an altitude of the trapezoid passing through point $B$ onto $AD$ at point $J$. Now, we can use the pythagorean theorem: $(22^2-(r-10)^2)+10^2=r^2$. Expanding and combining like terms gives us the quadratic 
 \[r^2-10r-242=0\] and solving for $r$ gives $r=5+\sqrt{267}$. So the solution is $5+267=\boxed{272}$.

 -The Geoman

 





\textbf{Solution 3}

Join the diameter of the circle $AD$ and let the length be $d$. By \href{https://artofproblemsolving.com/wiki/index.php?title=Ptolemy\%27s\_Theorem}{Ptolemy's Theorem} on trapezoid $ADEF$, $(AD)(EF) + (AF)(DE) = (AE)(DF)$. Since it is an isosceles trapezoid, both diagonals are equal. Let them be equal to $x$ each. Then

 
 \[20d + 22^2 = x^2\]

 Since $\angle AED$ is subtended by the diameter, it is right. Hence by the Pythagorean Theorem with right $\triangle AED$:

 
 \[(AE)^2 + (ED)^2 = (AD)^2\]
 \[x^2 + 22^2 = d^2\]

 From the above equations, we have:
 \[x^2 = d^2 - 22^2 = 20d + 22^2\]
 \[d^2 - 20d = 2\times22^2\]
 \[d^2 - 20d + 100 = 968+100 = 1068\]
 \[(d-10) = \sqrt{1068}\]
 \[d = \sqrt{1068} + 10 = 2\times(\sqrt{267}+5)\]

 Since the radius is half the diameter, it is $\sqrt{267}+5$, so the answer is $5+267 \Rightarrow \boxed{272}$.

 





\textbf{Solution 4}

As we can see this image, it is symmetrical hence the diameter divides the hexagon into two congruent quadrilateral. Now we can apply the Ptolemy's theorem. Denote the radius is r, we can get 
 \[22*2x+440=\sqrt{4x^2-400}\sqrt{4x^2-484}\], after simple factorization, we can get 
 \[x^4-342x^2-2420x=0\], it is easy to see that $x=-10, x=0$ are two solutions for the equation, so we can factorize that into 
 \[x(x+10)(x^2-10x-242)\]so we only need to find the solution for 
 \[x^2-10x-242=0\] and we can get $x=(\sqrt{267}+5)$ is the desired answer for the problem, and our answer is $5+267 \Rightarrow \boxed{272}$.～bluesoul

 





\textbf{Solution 6 (Trig Bash)}

Let $\angle{AOB} = \theta$. So, we have $\sin \dfrac{\theta}{2} = \dfrac{11}{r}$ and $\cos \dfrac{\theta}{2} = \dfrac{\sqrt{r^{2} - 121}}{r}$. So, $\sin \theta = 2 \sin \dfrac{\theta}{2} \cos \dfrac{\theta}{2} = \dfrac{22 \sqrt{r^{2} - 121}}{r^{2}}$. Let $H$ be the foot of the perpendicular from $B$ to $\overline{AD}$. We have $BF = 2 BH = 2 r \sin \theta = \dfrac{44 \sqrt{r^{2} - 121}}{r}$. Using Pythagorean theorem on $\triangle BCF$, to get $(\dfrac{44 \sqrt{r^{2} - 121}}{r})^{2} + 20^{2} = (2r)^{2}$, or $\dfrac{44^{2}r^{2} - 44^{2} \cdot 121}{r^{2}} + 20^{2} = 4r^{4}$. Multiplying by $r^{2}$, we get $44^{2} r^{2} - 44^{2} \cdot 121 + 20^{2} r^{2} = 4r^{4}$. Rearranging and simplifying, we get a quadratic in $r^{2}$: 
 \[r^{4} - 584r^{2} + 242^{2} = 0 \text{,}\] which gives us $r^{2} = 292 \pm 10\sqrt{267}$. Because $r$ is in the form $p + \sqrt{q}$, we know to choose the larger option, meaning $r^2 = 292 + 10\sqrt{267}$, so $p\sqrt{q} = 5\sqrt{267}$ and $p^2 + q = 292$. By inspection, we get $(p, q) = (5, 267)$, so our answer is $5 + 267 = \boxed{272}$.

 





\textbf{Solution 7}


                    \begin{figure}[H]
                        \centering
                        \includegraphics[width=0.5\textwidth]{Solution/AIME_2013_II_8_Solution_Figure_1.png}
                    \end{figure}
                    

 We know that $AD=x$ is a diameter, hence $ABD$ and $ACD$ are right triangles. Let $AB=BC=22$, and $CD=20.$ Hence, $ABD$ is a right triangle with legs $22,\sqrt{x^2-484},$ and hypotenuse, $x,$ and $ACD$ is a right triangle with legs $20, \sqrt{x^2-400},$ with hypotenuse $x$. By Ptolemy's we have 
 \[22(x+20)=\sqrt{x^2-400}\sqrt{x^2-484}\].We square both sides to get 
 \[484(x+20)^2=(x^2-400)(x^2-484) \implies 484(x+20)=(x-20)(x^2-484) \implies 484x=x^3-20x^2-484x \implies x(x^2-20x-968)=0\]

 We solve for $x$ via the Quadratic Formula and receive $x=10+2\sqrt{267}$, but we must divide by $2$ since we want the radius, and hence $267+5=\boxed{272}.$~SirAppel

 




	\item (\textit{AIME II 2013}) A $7\times 1$ board is completely covered by $m\times 1$ tiles without overlap; each tile may cover any number of consecutive squares, and each tile lies completely on the board. Each tile is either red, blue, or green. Let $N$ be the number of tilings of the $7\times 1$ board in which all three colors are used at least once. For example, a $1\times 1$ red tile followed by a $2\times 1$ green tile, a $1\times 1$ green tile, a $2\times 1$ blue tile, and a $1\times 1$ green tile is a valid tiling. Note that if the $2\times 1$ blue tile is replaced by two $1\times 1$ blue tiles, this results in a different tiling. Find the remainder when $N$ is divided by $1000$.


 
 \textbf{Answer: }106



\textbf{Solution 1}

Firstly, we consider how many different ways possible to divide the $7\times 1$ board.We ignore the cases of 1 or 2 pieces since we need at least one tile of each color.

 
                \begin{itemize}
                    \item Three pieces: $5+1+1$, $4+2+1$, $4+1+2$, etc, $\dbinom{6}{2}=15$ ways in total (just apply stars and bars here)\item Four pieces: $\dbinom{6}{3}=20$\item Five pieces: $\dbinom{6}{4}=15$\item Six pieces: $\dbinom{6}{5}=6$\item Seven pieces: $\dbinom{6}{6}=1$
                \end{itemize}


                Secondly, we use Principle of Inclusion-Exclusion to consider how many ways to color them:

 
                \begin{itemize}
                    \item Three pieces: $3^3-3\times 2^3+3=6$\item Four pieces: $3^4-3\times 2^4+3=36$\item Five pieces: $3^5-3\times 2^5+3=150$\item Six pieces: $3^6-3\times 2^6+3=540$\item Seven pieces: $3^7-3\times 2^7+3=1806$
                \end{itemize}


                Finally, we combine them together: $15\times 6+20\times 36+15\times 150+6\times 540+1\times 1806= 8106$.

 So the answer is $\boxed{106}$.

 \textbf{Explanation} In the $7*1$ tiling, when counting the ways to split this into $k$ pieces what you do is set the leftmost of the first tile and the rightmost of the last tile as borders.

 Then if we select say $1$ vertical line from the remaining $6$ vertical lines, we would have split into $2$ pieces. Similarly if we select $k$ vertical lines from them $6$, we would have $k+1$ pieces.

 Furthermore, to count the number of countings of $k$ pieces you apply the Principle of Inclusion and Exclusion.

 Done so by noticing that we have $3^k$ ways to color without any restricitons. Then, let's count the number of ways that we use atleast only $2$ colors and subtract off.

 If we use atleast only $2$ colors, we have $3$ ways to pick these colors, and then $3*2^k$ ways to color the squares in. However, this overcounts the case where we used only one color, two times each.

 Thus number of ways we use atleast only $2$ colors is $3*2^k-3$.

 Subtracting from total arrangements, we obtain $3^k-3*2^k+3$.

 





\textbf{Solution 1.1 (Faster/Streamlined 1)}

This solution is basically solution 1 with more things done at once. The game plan:

 $\sum_{i=0}^{7} ($the amount of ways to divide the board into $i$ pieces$) \cdot ($the amount of ways to color the respective divisions)

 The amount of ways to divide the board is determined using stars and bars. The colorings are found using PIE giving $3^i-3\cdot2^i+3$. Plus, we don't have to worry about the cases where $i=1$ or $i=2$ since they both give no solutions. So our equation becomes:

 $\sum_{i=3}^{7} \left(\dbinom{6}{i}\right)\cdot\left(3^i-3\cdot2^i+3\right)$

 Writing it all out and keeping the numbers small with mod 1000, we will eventually arrive at the answer of $\boxed{106}$.

 -NOTE: you can evaluate the above sum by binomial theorem: $(3+1)^6$ ,etc

 





\textbf{Solution 2 (Recursive)}

3 colors might be hard to calculate. How many ways can we tile an $n \times 1$ board with one color? It's going to be $2^{n-1}$ because if you draw out the $n$ spaces, you can decide whether each of the borders between the tiles are either there or not there. As there are $n-1$ borders, there are $2^{n-1}$ tilings for colorings with one color. Define a one-tiling of an $m \times 1$ as $f_1(m)$

 Now let's look at two colors. Let's see if we can get a two-tiling of an $(n+1) \times 1$ based off a $n \times 1$. There are $2$ cases to consider:

 1) The $n \times 1$ was a one-coloring and the block we are going to add consists of the second color. The number of ways we can do this is $2f_1(n)$

 2) The $n \times 1$ was a two-color tiling so now we've got $3$ cases to form the $(n+1) \times 1$: We can either add a block of the first color, the second color, or we can adjoin a block to the last block in the $n \times 1$. 

 This gives us $f_2(n+1)=2f_1(n)+3f_2(n)$

 Time to solve for the $3$-color tiling. We split into $2$ cases:

 1) The $n \times 1$ was a two-color tiling, and the block we are adding is of the 3rd color. This gives $f_2(n)$ ways but we have to multiply by $3 \choose 2$ $= 3$ because we have to pick 2 different colors for the two-color tiling. 

 2) The $n \times 1$ was a 3-color tiling, and we have to consider what we can do with the block that we are adding. It can be any of the 3 colors, or we can adjoin it to whatever was the last block in the $n \times 1$. This gives $4f_3(n)$

 So in total we have $f_3(n+1)=3f_2(n)+4f_3(n)$. 

 As shown by the below table:

 




	\item (\textit{AIME II 2013}) Given a circle of radius $\sqrt{13}$, let $A$ be a point at a distance $4 + \sqrt{13}$ from the center $O$ of the circle. Let $B$ be the point on the circle nearest to point $A$. A line passing through the point $A$ intersects the circle at points $K$ and $L$. The maximum possible area for $\triangle BKL$ can be written in the form $\frac{a - b\sqrt{c}}{d}$, where $a$, $b$, $c$, and $d$ are positive integers, $a$ and $d$ are relatively prime, and $c$ is not divisible by the square of any prime. Find $a+b+c+d$.


 
 \textbf{Answer: }146



\textbf{Solution 1 (Coordbash)}


                    \begin{figure}[H]
                        \centering
                        \includegraphics[width=0.5\textwidth]{Solution/AIME_2013_II_10_Solution_Figure_1.png}
                    \end{figure}
                    

 
Now we put the figure in the Cartesian plane, let the center of the circle $O (0,0)$, then $B (\sqrt{13},0)$, and $A(4+\sqrt{13},0)$

 The equation for Circle O is $x^2+y^2=13$, and let the slope of the line $AKL$ be $k$, then the equation for line $AKL$ is $y=k(x-4-\sqrt{13})$.

 Then we get $(k^2+1)x^2-2k^2(4+\sqrt{13})x+k^2\cdot (4+\sqrt{13})^2-13=0$. According to \href{https://artofproblemsolving.com/wiki/index.php?title=Vieta\%27s\_Formulas}{Vieta's Formulas}, we get

 $x_1+x_2=\frac{2k^2(4+\sqrt{13})}{k^2+1}$, and $x_1x_2=\frac{(4+\sqrt{13})^2\cdot k^2-13}{k^2+1}$

 So, $LK=\sqrt{1+k^2}\cdot \sqrt{(x_1+x_2)^2-4x_1x_2}$

 Also, the distance between $B$ and $LK$ is $\frac{k\times \sqrt{13}-(4+\sqrt{13})\cdot k}{\sqrt{1+k^2}}=\frac{-4k}{\sqrt{1+k^2}}$

 So the area $S=0.5ah=\frac{-4k\sqrt{(16+8\sqrt{13})k^2-13}}{k^2+1}$

 Then the maximum value of $S$ is $\frac{104-26\sqrt{13}}{3}$

 So the answer is $104+26+13+3=\boxed{146}$.

 





\textbf{Solution 2}


                    \begin{figure}[H]
                        \centering
                        \includegraphics[width=0.5\textwidth]{Solution/AIME_2013_II_10_Solution_Figure_1.png}
                    \end{figure}
                    

 Draw $OC$ perpendicular to $KL$ at $C$. Draw $BD$ perpendicular to $KL$ at $D$. 

 
 \[\frac{\triangle OKL}{\triangle BKL} = \frac{OC}{BD} = \frac{AO}{AB} = \frac{4+\sqrt{13}}{4}\]

 Therefore, to maximize area of $\triangle BKL$, we need to maximize area of $\triangle OKL$.

 
 \[\triangle OKL = \frac12 r^2 \sin{\angle KOL}\]

 So when area of $\triangle OKL$ is maximized, $\angle KOL = \frac{\pi}{2}$.

 Eventually, we get 
 \[\triangle BKL=  \frac12 \cdot (\sqrt{13})^2\cdot(\frac{4}{4+\sqrt{13}})=\frac{104-26\sqrt{13}}{3}\]

 So the answer is $104+26+13+3=\boxed{146}$.

 





\textbf{Solution 3 (simpler solution)}

A rather easier solution is presented in the Girls' Angle WordPress:

 \href{https://artofproblemsolving.comhttp://girlsangle.wordpress.com/2013/11/26/2013-aime-2-problem-10/}{http://girlsangle.wordpress.com/2013/11/26/2013-aime-2-problem-10/}

 


 





\textbf{Solution 4}

Let $N,M$ les on $AL$ such that $BM\bot AL, ON\bot AL$, call $BM=h, ON=k,LN=KN=d$ We call $\angle{LON}=\alpha$ By similar triangle, we have $\frac{h}{k}=\frac{4}{4+\sqrt{13}}, h=\frac{4k}{4+\sqrt{13}}$. Then, we realize the area is just $dh=d\cdot \frac{4K}{4+\sqrt{13}}$ As $\sin \alpha=\frac{d}{\sqrt{13}}, \cos \alpha=\frac{k}{\sqrt{13}}$. Now, we have to maximize $\frac{52\sin \alpha \cos \alpha}{4+\sqrt{13}}=\frac{26\sin 2\alpha}{4+\sqrt{13}}$, which is obviously reached when $\alpha=45^{\circ}$, the answer is $\frac{104-26\sqrt{13}}{3}$ leads to $\boxed{146}$

 





\textbf{Solution 5}


                    \begin{figure}[H]
                        \centering
                        \includegraphics[width=0.5\textwidth]{Solution/AIME_2013_II_10_Solution_Figure_1.png}
                    \end{figure}
                    Let C and D be the base of perpendiculars dropped from points O and B to AK.  Denote $BD = h, OC = H$.
 \[\triangle ABD \sim \triangle AOC \implies \frac {h}{H} = \frac {4}{4 + \sqrt{13}}.\]$KL$ is the base of  triangles $\triangle OKL$ and $\triangle BKL \implies \frac {[BKL]}{[OKL]} = \frac{h}{H} =$ const $\implies$The maximum possible area for $\triangle BKL$ and $\triangle OKL$ are at the same position of point $K$.

 $\triangle OKL$ has sides $OK = OL = \sqrt{13}\implies \max[\triangle OKL] = \frac {OK^2}{2} = \frac {13}{2}$

 in the case $\angle KOL = 90^\circ.$ It is possible – if we rotate such triangle, we can find position when $A$ lies on $KL.$ 
 \[\max[\triangle BKL] = \max[\triangle OKL] \cdot \frac {4}{4+\sqrt{13}} = \frac {26}{4+\sqrt{13}} \implies \boxed{\textbf{146}}\]\textbf{vladimir.shelomovskii@gmail.com, vvsss}

 




	\item (\textit{AIME II 2013}) Let $A = \{1, 2, 3, 4, 5, 6, 7\}$, and let $N$ be the number of functions $f$ from set $A$ to set $A$ such that $f(f(x))$ is a constant function. Find the remainder when $N$ is divided by $1000$.


 
 \textbf{Answer: }399



\textbf{Solution 1}

Any such function can be constructed by distributing the elements of $A$ on three tiers.

 The bottom tier contains the constant value, $c=f(f(x))$ for any $x$. (Obviously $f(c)=c$.)

 The middle tier contains $k$ elements $x\ne c$ such that $f(x)=c$, where $1\le k\le 6$.

 The top tier contains $6-k$ elements such that $f(x)$ equals an element on the middle tier.

 There are $7$ choices for $c$. Then for a given $k$, there are $\tbinom6k$ ways to choose the elements on the middle tier, and then $k^{6-k}$ ways to draw arrows down from elements on the top tier to elements on the middle tier.

 Thus $N=7\cdot\sum_{k=1}^6\tbinom6k\cdot k^{6-k}=7399$, giving the answer $\boxed{399}$.

 





\textbf{Solution 1 Clarified}

Define the three layers as \href{https://artofproblemsolving.com/wiki/index.php?title=Domain}{domain} $x$, \href{https://artofproblemsolving.com/wiki/index.php?title=Codomain}{codomain} $f(x)$, and codomain $f(f(x))$. Each one of them is contained in the \href{https://artofproblemsolving.com/wiki/index.php?title=Set}{set} $A$. We know that $f(f(x))$ is a \href{https://artofproblemsolving.com/wiki/index.php?title=Constant\_function}{constant function}, or in other words, can only take on one value. So, we can start off by choosing that value $c$ in $7$ ways. So now, we choose the values that can be $f(x)$ for all those values should satisfy $f(f(x))=c$. Let $S$ be that set of values. First things first, we must have $c$ to be part of $S$, for the $S$ is part of the domain of $x$. Since the values in $i\in S$ all satisfy $f(i) = c$, we have $c$ to be a value that $f(x)$ can be. Now, for the elements other than $5$:

 If we have $k$ elements other than $5$ that can be part of $S$, we will have $\binom{6}{k}$ ways to choose those values. There will also be $k$ ways for each of the elements in $A$ other than $c$ and those in set $S$ (for when \href{https://artofproblemsolving.com/wiki/index.php?title=Function}{function} $f$ is applied on those values, we already know it would be $c$). There are $6-k$ elements in $A$ other than $c$ and those in set $S$. Thus, there should be $k^{6-k}$ ways to match the domain $x$ to the values of $f(x)$. Summing up all possible values of $k$ ($[1,6]$), we have

 
 \[\sum_{k=1}^6 \binom{6}{k} k^{6-k} = 6\cdot 1 + 15\cdot 16 + 20\cdot 27 + 15\cdot 16 + 6\cdot 5 + 1 = 1057\]

 Multiplying that by the original $7$ for the choice of $c$, we have $7 \cdot 1057 = 7\boxed{399}.$

 


 





\textbf{Solution 2}

It is clear that we must have one fixed point (that is, $f(x)=x$). WLOG, let $x=1$ be a fixed point, so $f(1)=1$. 

 Now, let's do casework on how many of the inputs $2, 3, 4, 5, 6 ,7$ leads to $1$. Generally, if some values in that aforementioned list leads to $1$, then running it in the function again will yield $1$. All other values must be the the values that leads to $1$. 

 
For example:

 
$\textbf{Case 1:}$ All $6$ of $2, 3, 4, 5, 6, 7$ lead to $1$.In this case, there is only $1$ way.

 $\textbf{Case 2:}$ $5$ of $6$ of $2, 3, 4, 5, 6, 7$ lead to $1$. In this case, we choose $5$ of the $6$ to lead to $1$: $6\choose5$.

 Then, the other value that does not lead to $1$ should be one of the values that do: $5$ ways.

 $\binom{6}{5}\cdot5$

 $\textbf{Case 3:}$ $4$ of $6$ lead to $1$. Choose which lead to $1$: $6\choose4$ then the other values: $4^2$ ways

 $\binom{6}{4}\cdot4^2$

 $\textbf{Case 4:}$ $3$ of $6$ lead to $1$. $\binom{6}{3}\cdot3^3$

 $\textbf{Case 5:}$ $2$ of $6$ lead to $1$. $\binom{6}{2}\cdot2^4$

 $\textbf{Case 6:}$ $1$ of $6$ lead to $1$. $\binom{6}{1}\cdot1^5$

 
Adding up all the cases, we have $1057$ cases. But don't forget to account for the WLOG and multiply by $7$, yielding us the final answer of $7\boxed{399}.$

 


 





\textbf{Solution 3 (casework on range)}

Let $k$ be the element of set $A$ such that $f(f(a))=k$ for all $a \in A$. There are $7$ ways to choose $k$, and without loss of generality we can choose $k=1$.

 Denote by $R$ the range of $f$. For any $r \in R$, there exists $a \in A$ such that $f(a)=r$ so $f(f(a))=f(r)=1$. 

 Now we casework on the size of $R$.

 $\textbf{Case 1:} |R|=1$.

 The only element in $R$ is $1$, so there is only $1$ function in this case ($f(x)=1$ for all $x$).

 $\textbf{Case 2:} |R|=2$.

 There are $6$ ways to choose the other element of $R$. WLOG, let this element be $2$. We have $f(1)=f(2)=1$. For the other $5$ elements $e \notin A$, there are $2$ ways to determine $f(e)$ for every $e$, so there are $2^5$ functions $f$. However, the function $f(x) \equiv 1$ is excluded, as $2$ is in the range but there exists no $a$ such that $f(a)=2$. Therefore, there are $6 \cdot (2^5-1)=186$ functions in this case.

 $\textbf{Case 3:} |R|=3$.

 There are $\binom{6}{2}$ ways to choose the other $2$ elements of $R$. WLOG, let these elements be $2$ and $3$. We have $f(1)=f(2)=f(3)=1$. For the other $4$ elements $e \notin A$, there are $3$ ways to determine $f(e)$ for every $e$, so there are $3^4$ functions $f$. However, the $2^4$ functions that do not include $2$ in its range are excluded, and the $2^4$ functions that do not include $3$ in its range are also excluded. However, we have eliminated the function that excludes both $2$ and $3$ from its range, namely $f(x) \equiv 1$, twice, so we must add $1$ back. Therefore, there are $\binom{6}{2}(3^4-2^4-2^4+1)=750$ functions in this case.

 $\textbf{Case 4:} |R|=4$.

 There are $\binom{6}{3}$ ways to choose the other $3$ elements of $R$. WLOG, let these elements be $2, 3, 4$. We have $f(1)=f(2)=f(3)=f(4)=1$. For $e \in \{5, 6, 7\}$, if $f(e)=1$, then at least one of $2, 3, 4$ will not be included in the range of $f$. Therefore, each of $5, 6, 7$ must be mapped to a distinct element of $2, 3, 4$, so there are $3!$ ways to achieve this. Finally, the number of functions in this case is $\binom{6}{3}(3!)=120$.

 Finally, we add up all the cases to get $1057$ functions. However, this assumes $1$ is the constant value, but there are $7$ different possible constant values. Therefore, the final answer is $1057 \cdot 7 = 7399$ which when divided by $1000$ gives a remainder of $\boxed{399}$.

 


 




	\item (\textit{AIME II 2013}) Let $S$ be the set of all polynomials of the form $z^3 + az^2 + bz + c$, where $a$, $b$, and $c$ are integers. Find the number of polynomials in $S$ such that each of its roots $z$ satisfies either $|z| = 20$ or $|z| = 13$.


 
 \textbf{Answer: }540



\textbf{Solution 1}

Every cubic with real coefficients has to have either three real roots or one real and two nonreal roots which are conjugates. This follows from \href{https://artofproblemsolving.com/wiki/index.php?title=Vieta\%27s\_formulas}{Vieta's formulas}. 

 
                \begin{itemize}
                    \item Case 1:  $f(z)=(z-r)(z-\omega)(z-\omega^*)$, where $r\in \mathbb{R}$,  $\omega$ is nonreal, and $\omega^*$ is the complex conjugate of omega (note that we may assume that $\Im(\omega)>0$).
                \end{itemize}


                The real root $r$ must be one of $-20$, $20$, $-13$, or $13$. By Viète's formulas, $a=-(r+\omega+\omega^*)$, $b=|\omega|^2+r(\omega+\omega^*)$, and $c=-r|\omega|^2$. But $\omega+\omega^*=2\Re{(\omega)}$ (i.e., adding the conjugates cancels the imaginary part). Therefore, to make $a$ an integer, $2\Re{(\omega)}$ must be an integer. Conversely, if $\omega+\omega^*=2\Re{(\omega)}$ is an integer, then $a,b,$ and $c$ are clearly integers. Therefore $2\Re{(\omega)}\in \mathbb{Z}$ is equivalent to the desired property. Let $\omega=\alpha+i\beta$.

 
                \begin{itemize}
                    \item Subcase 1.1: $|\omega|=20$.
                \end{itemize}


                In this case, $\omega$ lies on a circle of radius $20$ in the complex plane. As $\omega$ is nonreal, we see that $\beta\ne 0$.  Hence $-20<\Re{(\omega)}< 20$, or rather $-40<2\Re{(\omega)}< 40$. We count $79$ integers in this interval, each of which corresponds to a unique complex number on the circle of radius $20$ with positive imaginary part.

 
                \begin{itemize}
                    \item Subcase 1.2: $|\omega|=13$.
                \end{itemize}


                In this case, $\omega$ lies on a circle of radius $13$ in the complex plane. As $\omega$ is nonreal, we see that $\beta\ne 0$.  Hence $-13<\Re{(\omega)}< 13$, or rather $-26<2\Re{(\omega)}< 26$. We count $51$ integers in this interval, each of which corresponds to a unique complex number on the circle of radius $13$ with positive imaginary part.

 Therefore, there are $79+51=130$ choices for $\omega$. We also have $4$ choices for $r$, hence there are $4\cdot 130=520$ total polynomials in this case.

 
                \begin{itemize}
                    \item Case 2: $f(z)=(z-r_1)(z-r_2)(z-r_3)$, where $r_1,r_2,r_3$ are all real.
                \end{itemize}


                In this case, there are four possible real roots, namely $\pm 13, \pm20$. Let $p$ be the number of times that $13$ appears among $r_1,r_2,r_3$, and define $q,r,s$ similarly for $-13,20$, and $-20$, respectively. Then $p+q+r+s=3$ because there are three roots. We wish to find the number of ways to choose nonnegative integers $p,q,r,s$ that satisfy that equation. By balls and urns, these can be chosen in $\binom{6}{3}=20$ ways.

 Therefore, there are a total of $520+20=\boxed{540}$ polynomials with the desired property.

 





\textbf{Solution 2 (Systematics)}

This combinatorics problem involves counting, and casework is most appropriate.There are two cases: either all three roots are real, or one is real and there are two imaginary roots.

 Case 1: Three roots are of the set ${13, -13, 20, -20}$. By stars and bars, there is $\binom{6}{3}=20$ ways (3 bars between all four possibilities, and then 3 stars that represent the roots themselves).

 Case 2: One real root: one of $13, -13, 20, -20$. Then two imaginary roots left; it is well known that because coefficients of the polynomial are integral (and thus not imaginary), these roots are conjugates. Therefore, either both roots have a norm (also called magnitude) of $20$ or $13$. Call the root $a+bi$, where $a$ is not the magnitude of the root; otherwise, it would be case 1. We need integral coefficients: expansion of $(x-(a+bi))(x-(a-bi))=-2ax+x^2+(a^2+b^2)$ tells us that we just need $2a$ to be integral, because $a^2+b^2$ IS the norm of the root! (Note that it is not necessary to multiply by the real root. That won't affect whether or not a coefficient is imaginary.) Therefore, when the norm is $20$, the $a$ term can range from $-19.5, -19, ...., 0, 0.5, ..., 19.5$ or $79$ solutions. When the norm is $13$, the $a$ term has $51$ possibilities from $-12.5, -12, ..., 12.5$. In total that's 130 total ways to choose the imaginary root. Now, multiply by the ways to choose the real root, $4$, and you get $520$ for this case.

 And $520+20=540$ and we are done.

 





\textbf{Solution 3 (Comments)}

If the polynomial has one real root and two complex roots, then it can be factored as $(z-r)(z^2+pz+q),$ where $r$ is real with $|r|=13,20$ and $p,q$ are integers with $p^2 <4q.$ The roots $z_1$ and $z_2$ are conjugates. We have $|z_1|^2=|z_2|^2=z_1z_2=q.$ So $q$ is either $20^2$ or $13^2$. The only requirement for $p$ is $p<\sqrt{4q^2}=2q.$ All such quadratic equations are listed as follows:

 $z^2+pz+20^2,$ where $p=0,\pm1,\pm2,\cdots,\pm 39,$

 $z^2+pz+13^2,$ where $p=0,\pm1,\pm2,\cdots,\pm 25$.

 Total of 130 equations, multiplied by 4 (the number of cases for real $r$, we have 520 equations, as indicated in the solution.

 -JZ

 





\textbf{Solution 4}

There are two cases: either all the roots are real, or one is real and two are imaginary.

 \textbf{Case 1:} All roots are real.Then each of the roots is a member of the set $\{-20, 20, -13, 13\}$. It splits into three sub-cases: either no two are the same, exactly two are the same, or all three are the same.

 \textbf{Sub-case 1.1:} No two are the same.This is obviously $\dbinom{4}{3}=4$.

 \textbf{Sub-case 1.2:} Exactly two are the same.There are four ways to choose the root that will repeat twice, and three ways to choose the remaining root. For this sub-case, $4\cdot 3=12$.

 \textbf{Sub-case 1.3:} All three are the same.This is obviously $4$.

 Thus for case one, we have $4+12+4=20$ polynomials in $S$. We now have case two, which we state below.

 \textbf{Case 2:} Two roots are imaginary and one is real.Let these roots be $p-qi$, $p+qi$, and $r$. Then by Vieta's formulas

 
                \begin{itemize}
                    \item $-(2p+r)=a$;\item $p^{2}+q^{2}+2pr=b$;\item $-\left(p^{2}+q^{2}\right)r=c$.
                \end{itemize}


                Since $a$, $b$, $c$, and $r$ are integers, we have that $p=\frac{1}{2}k$ for some integer $k$. Case two splits into two sub-cases now:

 \textbf{Sub-case 2.1:} $|p-qi|=|p+qi|=13$.Obviously, $|p|<13$. The $51$ cases in which $p$ is either $0,\pm\frac{1}{2},\pm\frac{2}{2},\pm\frac{3}{2},\ldots,\pm\frac{25}{2}$ are acceptable. Each can pair with one value of $q$ and four values of $r$, adding $51\cdot 4=204$ polynomials to $S$.

 \textbf{Sub-case 2.2:} $|p-qi|=|p+qi|=20$.Obviously, $|p|<20$. Here, the $79$ cases in which $p$ is either $0,\pm\frac{1}{2},\pm\frac{2}{2},\pm\frac{3}{2},\ldots,\pm\frac{39}{2}$ are acceptable. Again, each can pair with a single value of $q$ as well as four values of $r$, adding $79\cdot 4=316$ polynomials to $S$.

 Thus for case two, $204+316=520$ polynomials are part of $S$.

 All in all, $20+204+316=\boxed{540}$ polynomials can call $S$ home.

 


 


 




	\item (\textit{AIME II 2013}) In $\triangle ABC$, $AC = BC$, and point $D$ is on $\overline{BC}$ so that $CD = 3\cdot BD$. Let $E$ be the midpoint of $\overline{AD}$. Given that $CE = \sqrt{7}$ and $BE = 3$, the area of $\triangle ABC$ can be expressed in the form $m\sqrt{n}$, where $m$ and $n$ are positive integers and $n$ is not divisible by the square of any prime. Find $m+n$.


 
 \textbf{Answer: }010



\textbf{Solution 1}

We can set $AE=ED=m$. Set $BD=k$, therefore $CD=3k, AC=4k$. Thereafter

 




	\item (\textit{AIME II 2013}) For positive integers $n$ and $k$, let $f(n, k)$ be the remainder when $n$ is divided by $k$, and for $n > 1$ let $F(n) = \max_{\substack{1\le k\le \frac{n}{2}}} f(n, k)$. Find the remainder when $\sum\limits_{n=20}^{100} F(n)$ is divided by $1000$.


 
 \textbf{Answer: }512



\textbf{Solution}

\textbf{The Pattern} We can find that 

 $20\equiv 6 \pmod{7}$

 $21\equiv 5 \pmod{8}$

 $22\equiv 6 \pmod{8}$

 $23\equiv 7 \pmod{8}$

 $24\equiv 6 \pmod{9}$

 $25\equiv 7 \pmod{9}$

 $26\equiv 8 \pmod{9}$

 Observing these and we can find that the reminders are in groups of three continuous integers, considering this is true, and we get

 $99\equiv 31 \pmod{34}$

 $100\equiv 32 \pmod{34}$

 So the sum is $6+3\times(6+...+31)+31+32=1512$,it is also 17+20+23+...+95, so the answer is $\boxed{512}$.By: Kris17

 \textbf{The Intuition} First, let's see what happens if we remove a restriction. Let's define $G(x)$ as

 $G(x):=\max_{\substack{1\le k}} f(n, k)$

 Now, if you set $k$ as any number greater than $n$, you get n, obviously the maximum possible. That's too much freedom; let's restrict it a bit. Hence $H(x)$ is defined as

 $H(x):=\max_{\substack{1\le k\le n}} f(n, k)$

 Now, after some thought, we find that if we set $k=\lfloor \frac{n}{2} \rfloor+1$ we get a remainder of $\lfloor \frac{n-1}{2} \rfloor$, the max possible. Once we have gotten this far, it is easy to see that the original equation, 

 $F(n) = \max_{\substack{1\le k\le \frac{n}{2}}} f(n, k)$

 has a solution with $k=\lfloor \frac{n}{3} \rfloor+1$.

 \textbf{The Proof} The solution presented above does not prove why $F(x)$ is found by dividing $x$ by $3$. Indeed, that is the case, as rigorously shown below.

 Consider the case where $x = 3k$. We shall prove that $F(x) = f(x, k+1)$.For all $x/2\ge n > k+1, x = 2n + q$, where $0\le q< n$. This is because $x < 3k + 3 < 3n$ and $x \ge 2n$. Also, as $n$ increases, $q$ decreases. Thus, $q = f(x, n) < f(x, k+1) = k - 2$ for all $n > k+1$.Consider all $n < k+1. f(x, k) = 0$ and $f(x, k-1) = 3$. Also, $0 < f(x, k-2) < k-2$. Thus, for $k > 5, f(x, k+1) > f(x, n)$ for $n < k+1$.

 Similar proofs apply for $x = 3k + 1$ and $x = 3k + 2$. The reader should feel free to derive these proofs themself.

 \textbf{Generalized Solution} $Lemma:$ Highest remainder when $n$ is divided by $1\leq k\leq n/2$ is obtained for $k_0 = (n + (3 - n \pmod{3}))/3$ and the remainder thus obtained is $(n - k_0*2) = [(n - 6)/3 + (2/3)*n \pmod{3}]$. 

 $Note:$ This is the second highest remainder when $n$ is divided by $1\leq k\leq n$ and the highest remainder occurs when $n$ is divided by $k_M$ = $(n+1)/2$ for odd $n$ and  $k_M$ = $(n+2)/2$ for even $n$.

 Using the lemma above:

 $\sum\limits_{n=20}^{100} F(n) = \sum\limits_{n=20}^{100} [(n - 6)/3 + (2/3)*n \pmod{3}]$$= [(120*81/2)/3 - 2*81 + (2/3)*81] = 1512$

 So the answer is $\boxed{512}$

 Proof of Lemma: It is similar to $The Proof$ stated above. 

 Kris17

 





\textbf{Solution 2}

Let $f(n,k)=r$, and we set the quotient $q=\left\lfloor\frac{n}{k}\right\rfloor\ge 2$. It follows that
 \[r=n-qk\]We see that
 \[q\le \frac{n}{k} < q+1\]Or
 \[\frac{n}{q+1}< k\le\frac{n}{q}\]Minimizing $k$ yields a greater value for $r$. Hence, taking $k=\left\lfloor\frac{n}{q+1}\right\rfloor+1$,
 \[r=n-q\left\lfloor\frac{n}{q+1}\right\rfloor-q\]When $q=2$:
 \[r=n-2\left\lfloor\frac{n}{3}\right\rfloor-2 \ge \frac{n}{3}-2\]Since $\left\lfloor\frac{n}{q+1}\right\rfloor\ge \frac{n}{q+1}-1$, comparing to the case when $q\ge 3$
 \[r\le \frac{n}{q+1}\le \frac{n}{4}\]Thus, the maximal of $r$ holds when $q=2,n\ge 24$. Checking the cases for $n=20,21,22,23$, $q=2$ also gives the maximal.

 For $n\ge 20$, we have
 \[F(n)=n-2\left\lfloor\frac{n}{3}\right\rfloor-2\]Computing the summation:
 \[\sum_{n=20}^{100}F(n)=\left(\sum_{n=20}^{100}n\right)-\left(2\sum_{n=20}^{100}\left\lfloor\frac{n}{3}\right\rfloor\right)-\sum_{n=20}^{100}2\]We see that\begin{align*}\sum_{n=20}^{100}\left\lfloor\frac{n}{3}\right\rfloor&=\sum_{k=7}^{33}\left\lfloor\frac{3k-1}{3}\right\rfloor+\sum_{k=7}^{33}\left\lfloor\frac{3k}{3}\right\rfloor+\sum_{k=7}^{33}\left\lfloor\frac{3k+1}{3}\right\rfloor\\&=\sum_{k=7}^{33}(k-1)+\sum_{k=7}^{33}k+\sum_{k=7}^{33}k\\&=\sum_{k=7}^{33}3k-1\\&=1593\end{align*}
Therefore,
 \[\sum_{n=20}^{100}F(n)=4860-3186-162=1512\]The requested answer is $\boxed{512}$.

 


 




	\item (\textit{AIME II 2013}) Let $A,B,C$ be angles of an acute triangle with

 \begin{align*} \cos^2 A + \cos^2 B + 2 \sin A \sin B \cos C &= \frac{15}{8} \text{ and} \\ \cos^2 B + \cos^2 C + 2 \sin B \sin C \cos A &= \frac{14}{9} \end{align*}

There are positive integers $p$, $q$, $r$, and $s$ for which 
 \[\cos^2 C + \cos^2 A + 2 \sin C \sin A \cos B = \frac{p-q\sqrt{r}}{s},\] where $p+q$ and $s$ are relatively prime and $r$ is not divisible by the square of any prime.  Find $p+q+r+s$.


 
 \textbf{Answer: }222



\textbf{Solution 1}

Let's draw the triangle. Since the problem only deals with angles, we can go ahead and set one of the sides to a convenient value. Let $BC = \sin{A}$.

 By the Law of Sines, we must have $CA = \sin{B}$ and $AB = \sin{C}$.

 Now let us analyze the given:

 
 \begin{align*} \cos^2A + \cos^2B + 2\sin A\sin B\cos C &= 1-\sin^2A + 1-\sin^2B + 2\sin A\sin B\cos C \\ &= 2-(\sin^2A + \sin^2B - 2\sin A\sin B\cos C) \end{align*}


 Now we can use the Law of Cosines to simplify this:

 
 \[= 2-\sin^2C\]

 
Therefore: 
 \[\sin C = \sqrt{\dfrac{1}{8}},\cos C = \sqrt{\dfrac{7}{8}}.\] Similarly, 
 \[\sin A = \sqrt{\dfrac{4}{9}},\cos A = \sqrt{\dfrac{5}{9}}.\] Note that the desired value is equivalent to $2-\sin^2B$, which is $2-\sin^2(A+C)$. All that remains is to use the sine addition formula and, after a few minor computations, we obtain a result of $\dfrac{111-4\sqrt{35}}{72}$. Thus, the answer is $111+4+35+72 = \boxed{222}$.

 Note that the problem has a flaw because $\cos B < 0$ which contradicts with the statement that it's an acute triangle. Would be more accurate to state that $A$ and $C$ are smaller than 90. -Mathdummy

 





\textbf{Solution 2}

Let us use the identity $\cos^2A+\cos^2B+\cos^2C+2\cos A \cos B \cos C=1$ .

 Add 
 \begin{align*} \cos^2 C+2(\cos A\cos B-\sin A \sin B)\cos C \end{align*}
 to both sides of the first given equation. 

 


 
Thus, as 
 \begin{align*} \cos A\cos B-\sin A\sin B=\cos (A+B)=-\cos C ,\end{align*}
 we have 
 \begin{align*} \dfrac{15}{8}-2\cos^2 C +\cos^2 C=1,  \end{align*}
 so $\cos C$ is $\sqrt{\dfrac{7}{8}}$ and therefore $\sin C$ is $\sqrt{\dfrac{1}{8}}$. 

 Similarily, we have $\sin A =\dfrac{2}{3}$ and $\cos A=\sqrt{\dfrac{14}{9}-1}=\sqrt{\dfrac{5}{9}}$ and the rest of the solution proceeds as above.

 





\textbf{Solution 3}

Let 
 \begin{align*} \cos^2 A + \cos^2 B + 2 \sin A \sin B \cos C &= \frac{15}{8} \text{     ------ (1)}\\ \cos^2 B + \cos^2 C + 2 \sin B \sin C \cos A &= \frac{14}{9} \text{     ------ (2)}\\ \cos^2 C + \cos^2 A + 2 \sin C \sin A \cos B &= x \text{                        ------ (3)}\\ \end{align*}


 Adding (1) and (3) we get:
 \[2 \cos^2 A + \cos^2 B + \cos^2 C + 2 \sin A( \sin B \cos C +  \sin C \cos B)  = \frac{15}{8} + x\]   or
 \[2 \cos^2 A + \cos^2 B + \cos^2 C + 2 \sin A \sin (B+C)  = \frac{15}{8} + x\]    or
 \[2 \cos^2 A + \cos^2 B + \cos^2 C + 2 \sin ^2 A = \frac{15}{8} + x\]    or
 \[\cos^2 B + \cos^2 C =  x - \frac{1}{8}     \text{     ------ (4)}\]

 Similarly adding (2) and (3) we get:
 \[\cos^2 A + \cos^2 B =  x - \frac{4}{9}     \text{     ------ (5)}\]Similarly adding (1) and (2) we get:
 \[\cos^2 A + \cos^2 C =  \frac{14}{9}  - \frac{1}{8}     \text{     ------ (6)}\]

 And (4) - (5) gives:
 \[\cos^2 C - \cos^2 A =  \frac{4}{9}  - \frac{1}{8}    \text{     ------ (7)}\]

 Now (6) - (7) gives:$\cos^2 A  =  \frac{5}{9}$ or$\cos A = \sqrt{\dfrac{5}{9}}$ and $\sin A = \frac{2}{3}$so $\cos C$ is $\sqrt{\dfrac{7}{8}}$ and therefore $\sin C$ is $\sqrt{\dfrac{1}{8}}$

 Now $\sin B = \sin(A+C)$ can be computed first and then $\cos^2 B$ is easily found.

 Thus $\cos^2 B$ and $\cos^2 C$ can be plugged into (4) above to give x =  $\dfrac{111-4\sqrt{35}}{72}$.

 Hence the answer is = $111+4+35+72 = \boxed{222}$.

 Kris17

 





\textbf{Solution 4}

Let's take the first equation $\cos^2 A + \cos^2 B + 2 \sin A \sin B \cos C = \frac{15}{8}$. Substituting $180 - A - B$ for C, given A, B, and C form a triangle, and that $\cos C = \cos(A + B)$, gives us:

 $\cos^2 A + \cos^2 B + 2 \sin A \sin B \cos (A+B) = \frac{15}{8}$

 Expanding out gives us $\cos^2 A + \cos^2 B + 2 \sin^2 A \sin^2 B - 2 \sin A \sin B \cos A \cos B = \frac{15}{8}$. 

 Using the double angle formula $\cos^2 k = \frac{\cos (2k) + 1}{2}$, we can substitute for each of the squares $\cos^2 A$ and $\cos^2 B$. Next we can use the Pythagorean identity on the $\sin^2 A$ and $\sin^2 B$ terms. Lastly we can use the sine double angle to simplify.

 $\cos^2 A + \cos^2 B + 2(1 - \cos^2 A)(1 - \cos^2 B) - \frac{1}{2} \cdot \sin 2A \sin 2B = \frac{15}{8}$. 

 Expanding and canceling yields, and again using double angle substitution, 

 $1 + 2 \cdot \frac{\cos (2A) + 1}{2} \cdot \frac{\cos (2B) + 1}{2} - \frac{1}{2} \cdot \sin 2A \sin 2B = \frac{15}{8}$. 

 Further simplifying yields:

 $\frac{3}{2} + \frac{\cos 2A \cos 2B - \sin 2A \sin 2B}{2} = \frac{15}{8}$. 

 Using cosine angle addition formula and simplifying further yields, and applying the same logic to Equation $2$ yields:

 $\cos (2A + 2B) = \frac{3}{4}$ and $\cos (2B + 2C) = \frac{1}{9}$.

 Substituting the identity $\cos (2A + 2B) = \cos(2C)$, we get:

 $\cos (2C) = \frac{3}{4}$ and $\cos (2A) = \frac{1}{9}$.

 Since the third expression simplifies to the expression $\frac{3}{2} + \frac{\cos (2A + 2C)}{2}$, taking inverse cosine and using the angles in angle addition formula yields the answer, $\frac{111 - 4\sqrt{35}}{72}$, giving us the answer $\boxed{222}$.

 





\textbf{Solution 5}

We will use the sum to product formula to simplify these equations. Recall 
 \[2\sin{\frac{\alpha+\beta}{2}}\sin{\frac{\alpha-\beta}{2}} = \cos{\alpha}+\cos{\beta}.\] Using this, let's rewrite the first equation: 
 \[\cos^2(A) + \cos^2(B) + 2 \sin(A) \sin(B) \cos(C) = \frac{15}{8}\] 
 \[\cos^2(A) + \cos^2(B) + (\cos(A+B)+\cos(A-B))\cos(C).\] Now, note that $\cos(C)=-\cos(A+B)$. 
 \[\cos^2(A) + \cos^2(B) + (\cos(A+B)+\cos(A-B))(-\cos(A+B))\] 
 \[\cos^2(A) + \cos^2(B) - \cos(A+B)\cos(A-B)+cos^2(A+B)=\frac{15}{8}.\] We apply the sum to product formula again. 
 \[\cos^2(A) + \cos^2(B) - \frac{\cos(2A)+\cos(2B)}{2}+cos^2(A+B)=\frac{15}{8}.\] Now, recall that $\cos(2\alpha)=2\cos^2(\alpha)-1$. We apply this and simplify our expression to get: 
 \[\cos^2(A+B)=\frac{7}{8}\] 
 \[\cos^2(C)=\frac{7}{8}.\] Analogously, 
 \[\cos^2(A)=\frac{5}{9}.\] 
 \[\cos^2(A+C)=\frac{p-q\sqrt{r}}{s}-1.\] We can find this value easily by angle sum formula. After a few calculations, we get $\frac{111 - 4\sqrt{35}}{72}$, giving us the answer $\boxed{222}$.~superagh

 





\textbf{Solution 6}

According to LOC $a^2+b^2-2ab\cos{\angle{c}}=c^2$, we can write it into $\sin^2{\angle{A}}+\sin^2{\angle{B}}-2\sin{\angle{A}}\sin{\angle{B}}\cos{\angle{C}}=\sin^2{\angle{C}}$. $\sin^2{\angle{A}}+\sin^2{\angle{B}}-2\sin{\angle{A}}\sin{\angle{B}}\cos{\angle{C}}+cos^2A+cos^2B+2sinAsinBcosC=\frac{15}{8}+sin^2C$ We can simplify to $2=sin^2C+\frac{15}{8}$. Similarly, we can generalize $2=sin^2A+\frac{14}{9}$. After solving, we can get that $sinA=\frac{2}{3}; cosA=\frac{\sqrt{5}}{3}; sinC=\frac{\sqrt{2}}{4}; cosC=\frac{\sqrt{14}}{4}$Assume the value we are looking for is $x$, we get $sin^2B+x=2$, while $sinB=sin(180^{\circ}-A-C)=sin(A+C)$ which is $\frac{2\sqrt{14}+\sqrt{10}}{12}$, so $x=\frac{111 - 4\sqrt{35}}{72}$, giving us the answer $\boxed{222}$.~bluesoul

 


 




\end{enumerate}